% Complete author proof developed from the physical source, not a referee report.
% Dependencies: native 260, 279, 281--283, 316--317, and new 319.
% No change to the section, true collision record, physical roof, or outer limits.
\section{Uniform total turning of the physical marked boundary}
\label{sec:ab-turning}

The boundary measure in Theorem~\ref{thm:bf-flux} cannot be bounded
pointwise by deleting its nonzero critical offset.  We instead control
its complete angular variation on the actual marked source.  The
argument uses the physical one-sided interval partitions of
Theorem~\ref{thm:mb-branch-sup} and their two-sided weighted overlay
in Theorem~\ref{thm:mb-two-sided-sum}.  In particular its partition
is not a hypothetical family of large comparison cells.

Keep the original endpoint cap of Lemma~\ref{lem:v109-section}:
\[
 |p_0|,|p_m|\le 9/10,\qquad c_\circ=\sqrt{19}/10.
\]
This cap follows from the original section tests.  Restricting an
upper-bound source to it adds no new condition to the observed event.
Let $k_-=100/47$, $k_+=20/9$, $U=170/9$,
$a_\circ=\sqrt{1827}/200$, $b_*=97/200$, and
$\varrho=47/59$.  Use exactly the fixed $B_{\rm cv}$, $k_0$,
$\ell_*$ and physical curve class of
Theorem~\ref{thm:uc-growth}.  Set
\begin{align}
 B_2&=B_{\rm cv}(1+U^2)^{3/2},&
 \Gamma&=\sqrt{1+k_+^2},&
 L_0&=\frac{\pi\sqrt{1+U^2}}{k_-},\notag\\
 L_h&=\frac{3\sqrt{1+U^2}}{k_-},&
 c_H&=\sin(k_0^{-2}),\notag\\
 V_H&=\max\left\{-2\log c_H,
            \log(\pi/2)+2\log(1+k_0^{-1})\right\}.
                                                        \label{eq:ab-constants-one}
\end{align}
Let $D_B$ be the completed one-step derivative constant in
\eqref{eq:uc-DB}, and $K_0$ the first-image derivative constant in
\eqref{eq:ag-geometry-constants}.  The following constants are finite
and depend only on the original physical family and its already fixed
curve class:
\begin{align}
 D_1^{\rm var}
  &=\max\left\{
       (U/c_H+D_B)L_0^{2/3},
       (2\pi U+D_B)L_h^{2/3}\right\},\notag\\
 V_J&=\frac{D_1^{\rm var}\Gamma^{1/3}L_0^{1/3}}
                     {1-\varrho^{1/3}},\notag\\
 V_\rho&=V_H+
       \left(a_\circ^{-1}+b_*a_\circ^{-2}+K_0/2\right)\ell_*,
                                                        \label{eq:ab-constants-two}\\
 V_t&=-2\log c_\circ+B_2L_0/k_-,\notag\\
 V_1&=\frac{L_0(k_+b_*+1)}{k_-a_\circ^2},\notag\\
 V_A&=\max\{V_1,V_\rho+V_J+V_t\}.
                                                        \label{eq:ab-constants-three}
\end{align}
No constant is tuned to the collision count, a clearance level, a
roof window, an exact lattice label, or an outer cutoff.

\begin{lemma}[Absolute homogeneous logarithmic variation]
\label{lem:ab-jacobian-variation}
Let $C'$ be a connected actual curve branch in the class of
Theorem~\ref{thm:uc-growth}, and suppose its images through time $l$
are homogeneous.  If $L_l$ is the arclength of its final image, then
\begin{equation}\label{eq:ab-jacobian-variation}
 \operatorname{Var}_{C'}\log J_C(T_R^l)
 \le \frac{D_1^{\rm var}\Gamma^{1/3}}
                  {1-\varrho^{1/3}} L_l^{1/3}
 \le V_J.
\end{equation}
The variation is along the full indicated connected branch, not just
the difference between its endpoint values.  No positive lower bound
on the interior incidences is imposed.
\end{lemma}
\begin{proof}
Every unstable image graph has $k_-\le d\varphi/ds\le U$ and
$-\pi/2<\varphi<\pi/2$.  Integration in $\varphi$ therefore
bounds its whole arclength by $L_0$.  On a segment of length $r$ in
$H_k^\pm$ one has
\[
 r\le L_hk^{-3},\qquad
 c\ge (2/\pi)(k+1)^{-2}.
\]
Lemma~\ref{lem:uc-logjac} bounds the absolute logarithmic derivative
per unit image arclength by $U/c+D_B$.  Thus its absolute integral
on this segment is at most
\[
 (2\pi U+D_B)k^2r
 \le (2\pi U+D_B)L_h^{2/3}r^{1/3}.
\]
On $H_0$ the analogous bound is
$(U/c_H+D_B)r\le(U/c_H+D_B)L_0^{2/3}r^{1/3}$.
These are estimates of the integral of the absolute derivative.
They prove a one-step \emph{total-variation} estimate with constant
$D_1^{\rm var}$; no conclusion about variation is inferred merely
from a H\"older difference bound.

The exact adapted expansion of \eqref{eq:uc-adapted} gives
$L_i\le\Gamma\varrho^{l-i}L_l$ for every intermediate image.
Write the logarithm of the full physical Jacobian as the sum of
its $l$ actual one-step logarithms.  Pullback by a monotone curve
parametrization preserves their total variations.  The triangle
inequality for variation and the preceding absolute one-step bounds
give
\[
 \operatorname{Var}_{C'}\log J_C(T_R^l)
 \le D_1^{\rm var}\sum_{i=1}^l L_i^{1/3}
 \le D_1^{\rm var}\Gamma^{1/3}L_l^{1/3}
                     \sum_{r=0}^{l-1}\varrho^{r/3}.
\]
Sum the convergent geometric series.  The $l=0$ Jacobian is one
and has zero variation.  Compact subarcs and monotone exhaustion
give the same inequalities on open homogeneous intervals, including
arbitrarily short intervals whose endpoints approach singularities.
\end{proof}

\begin{lemma}[Uniform variation of a physical half-ray momentum derivative]
\label{lem:ab-momentum-variation}
On every one-sided angular branch of
Theorem~\ref{thm:mb-branch-sup}, let
\[
 A_l^p(\vartheta)=|\partial_\vartheta p_l|>0
\]
be the actual momentum derivative after $l\ge1$ true collisions
from the marked foot.  Restrict the terminal momentum to the original
cap $|p_l|\le9/10$.  This restriction is an interval, possibly empty,
and on it
\begin{equation}\label{eq:ab-momentum-variation}
             \operatorname{Var}\log A_l^p\le V_A.
\end{equation}
The same assertion holds in the reversed physical direction, with
the transported original event and the same constants.  Neither a
terminal homogeneous strip count nor an extra collision is inserted.
\end{lemma}
\begin{proof}
The terminal curve is a genuine unstable graph in its direction of
propagation.  Its angle $\varphi_l$ and momentum $p_l$ are strictly
monotone in its curve coordinate.  The cap is consequently one interval.
Since $k_0\ge2$ and
$\arcsin(9/10)<\pi/2-1/4$, the cap is contained in the original
central homogeneous strip $H_0$.

For $l\ge2$, start from the exact first-image source
$\rho_C\,d\ell_C$ of Theorem~\ref{thm:ag-family}.  Its first image
$C$ is already homogeneous and has length at most $\ell_*$.  The
one-sided branch in Theorem~\ref{thm:mb-branch-sup} has homogeneous
intermediate images; its last true collision, restricted to the cap,
has its image in $H_0$.  The arclength chain rule gives the exact
identity
\begin{equation}\label{eq:ab-A-density}
 A_l^p
 =\frac{c_lV_l}{\sqrt{1+V_l^2}}
          \frac{J_C(T_R^{l-1})}{\rho_C}.
\end{equation}
Here $\rho_C$ is the actual angular density before any Borel record
selector is multiplied.  The selector is not differentiated.

We verify each variation in \eqref{eq:ab-A-density}.  On a first-image
homogeneous graph,
\[
 \log\rho_C=\log c_1-\log a-\tfrac12\log(1+V_1^2).
\]
The angle $\varphi_1$ is monotone.  On $H_0$ the variation of
$\log c_1$ is at most $-2\log c_H$; on a signed $H_k$ it is at
most
\[
 \log\!\frac{\sin(k^{-2})}{\sin((k+1)^{-2})}
 \le\log(\pi/2)+2\log(1+k_0^{-1}).
\]
Both cases are bounded by $V_H$.  The exact first-flight derivatives
in \eqref{eq:ag-a-derivative} yield
\[
 |a_s|\le1+b_*/a_\circ,
 \qquad |V_{1,s}|\le K_0.
\]
Using $a\ge a_\circ$ and $v/(1+v^2)\le1/2$ for $v>0$, their
absolute integrals on the first piece show
\[
                 \operatorname{Var}\log\rho_C\le V_\rho.
\]
This retains the true grazing factor in $\rho_C$.  It does not
replace the density by a constant before summing the initial strips.

Lemma~\ref{lem:ab-jacobian-variation} gives variation at most $V_J$
for the Jacobian in \eqref{eq:ab-A-density}.  At the terminal cap
$c_l\ge c_\circ$ and $\varphi_l$ is monotone, so
\[
 \operatorname{Var}\log c_l\le-2\log c_\circ.
\]
Geometric curvature gives $|V_{l,s}|\le B_2$.  Moreover
\[
 \left|\partial_s\log\frac{V_l}{\sqrt{1+V_l^2}}\right|
       =\frac{|V_{l,s}|}{V_l(1+V_l^2)}\le B_2/k_-.
\]
Its absolute integral is at most $B_2L_0/k_-$.  Thus the terminal
factor has logarithmic variation at most $V_t$.  Taking variation
of \eqref{eq:ab-A-density} proves the second alternative in
\eqref{eq:ab-constants-three}.

For $l=1$, the branch partition uses the finite first-hit intervals,
without an unnecessary infinite homogeneous cut.  Put
$b=\varsigma(R+u_0)$.  The exact first-hit derivative is
$A_1^p=\kappa a+c$, where $\kappa=1/R$.  Directly from
\eqref{eq:ag-a-derivative},
\[
 \partial_s(\kappa a+c)
 =\kappa(p+bc/a)-p(\kappa+c/a)
 =c(\kappa b-p)/a.
\]
Since $A_1^p\ge k_-a_\circ$, its logarithmic derivative has magnitude
at most $(k_+b_*+1)/(k_-a_\circ^2)$.  The first image has length
at most $L_0$, proving the bound $V_1$.

Time reversal is the exact arclength isometry already used in
Theorems~\ref{thm:uc-growth} and \ref{thm:nm-kernel}.  It reverses
the momentum sign and transports the original selector; it preserves
all the absolute derivative and variation bounds above.  The same
proof therefore covers the $j+1$ backward collisions, including
$j=0$.  All limiting statements use compact regular subarcs followed
by exhaustion, not a fictitious reflection at the marked foot.
\end{proof}

\begin{lemma}[Radial turning controlled by tangent turning]
\label{lem:ab-radial-turning}
Let $P$ be a $C^2$ regular plane curve on a connected interval.
Suppose $P$ is nonzero and its tangent has total turning $V_\psi$.
Then
\begin{equation}\label{eq:ab-radial-turning}
           \int |d\arg P|\le V_\psi+\pi.
\end{equation}
If $P$ passes through the origin at most once, the sum on its punctured
intervals is at most $V_\psi+2\pi$.  The latter statement excludes
the jump of the polar angle at the origin; it assigns no critical fan
there a zero value.
\end{lemma}
\begin{proof}
On a compact nonzero subarc choose continuous angles
$\alpha=\arg P$ and $\psi=\arg P'$.  With
$\delta=\psi-\alpha$, the exact derivative is
\[
         \alpha'=\frac{|P'|}{|P|}\sin\delta.
\]
The $2\pi$-periodic triangular function
$h(\delta)=\arccos(\cos\delta)$ takes values in $[0,\pi]$ and
has derivative $\operatorname{sgn}(\sin\delta)$ away from multiples
of $\pi$.  The absolutely continuous chain rule gives, almost
everywhere,
\[
 |\alpha'|
   =\operatorname{sgn}(\sin\delta)\psi'
                                -\frac{d}{ds}h(\delta).
\]
At a level set $\delta\in\pi\mathbb Z$, its derivative is zero
almost everywhere on that set; also $\alpha'=0$ there, so the same
identity holds with the sign set to zero.  Integration gives
\[
 \int|\alpha'|
 \le\int|\psi'|+h(\delta_{\rm left})-h(\delta_{\rm right})
 \le V_\psi+\pi.
\]
Exhausting the interval proves the first assertion.  If there is
one zero, apply it on both punctured sides and sum their tangent
variations.  This proves the second assertion and its explicit
critical-point convention.
\end{proof}

\begin{corollary}[Uniform turning on every actual two-sided angular piece]
\label{cor:ab-complete-piece-turning}
On any nonempty overlap piece from
Theorem~\ref{thm:mb-two-sided-sum}, restricted to the original
endpoint cap, set $P(\vartheta)=(-p_0(\vartheta),p_m(\vartheta))$.
Then
\begin{equation}\label{eq:ab-piece-turning}
 \int_{\{P\ne0\}}|\partial_\vartheta\arg P|\,d\vartheta
                         \le 2\pi+V_A.
\end{equation}
The bound is independent of the length of the overlap and of its
distance from a critical record.  All original Borel cuts may be
retained as selectors without counting their connected components.
\end{corollary}
\begin{proof}
The signed derivative identity \eqref{eq:bf-signed-derivatives}
makes both coordinate derivatives of $P$ nonzero with fixed signs
on the piece.  Its tangent therefore stays in one open quadrant,
and $P$ can pass through the origin at most once.  The cap restricts
each monotone momentum to an interval, so their simultaneous
restriction is still an interval.

The absolute slope ratio of the tangent is $A_+^p/A_-^p$.  On its
fixed quadrant the derivative of its tangent angle with respect to
$\log(A_+^p/A_-^p)$ has absolute value
\[
 \frac{A_+^p/A_-^p}{1+(A_+^p/A_-^p)^2}\le\frac12.
\]
Lemma~\ref{lem:ab-momentum-variation} bounds each of the two
logarithmic variations by $V_A$, also after restriction to the
overlap.  Hence the tangent turning is at most $V_A$.  Apply
Lemma~\ref{lem:ab-radial-turning}.  A Borel selector bounded by one
only decreases the nonnegative integral already bounded; it is
neither differentiated nor used to create additional summation pieces.
\end{proof}

\section{The complete regular clearance-boundary budget without an offset cutoff}
\label{sec:ab-all-offset}

Let $\mathfrak B^{\rm reg}_{j,u_0}$ be the original positive marked
charge of \eqref{eq:bf-flux}, for one specified incident side of the
clearance level $u_0\in[0,3/200]$, summed over both impact signs and
all actual regular angular branches of the $m$-flight record.
Here $0\le j<m$, all visited collisions on an integrated branch are
regular, and all original endpoint and exact-record tests are retained.
The regular boundary is punctured only at $g=0$, where the original
incoming critical fan remains a different measure.  There is no
positive lower bound on $|g|$ or on an interior incidence.

For $u_0=0$ this is the actual $u>0$ one-sided unvisited trace of
Theorem~\ref{thm:nm-trace}; no collision is added at the mark.
An isolated singular angular parameter is not used to define a
new trace.  In particular the notation does not absorb a distinct
visited-grazing, section or pairing boundary, or a simultaneous
critical corner fan.  Every regular part of the prescribed marked
boundary, including nonzero-offset and arbitrarily short components,
is included.

Put
\begin{equation}\label{eq:ab-final-constant}
 C_{\rm ab}
    =\frac{\Lambda_\circ(2\pi+V_A)}{4\pi c_*R_-^2}\,
       \frac{607}{90}\,H_{\rm nm}\,\frac{94}{\sqrt{1827}}.
\end{equation}
The last three factors are exactly the complete two-sided root-weight
constant in \eqref{eq:mb-two-sided-sum}; in particular $H_{\rm nm}$
already includes both impact signs.

\begin{theorem}[Full regular marked boundary, all critical offsets]
\label{thm:ab-all-offset}
For the unchanged Lorentz family and the complete indicated physical
source,
\begin{equation}\label{eq:ab-all-offset}
 \mathfrak B^{\rm reg}_{j,u_0}
       \le C_{\rm ab}
                  \bigl(\varrho^j+\varrho^{m-1-j}\bigr).
\end{equation}
This bound is uniform in $R$, $m$, the level $u_0$, and every original
exact record and roof selector between zero and one.  It requires no
critical zero-offset chart, no critical momentum cutoff, and no
whole-word pressure or comparison-pool abundance premise.

Consequently, for any levels $u_j\in[0,3/200]$ and any integer
$D\ge0$,
\begin{equation}\label{eq:ab-full-depth}
 \sum_{\min(j,m-1-j)\ge D}\mathfrak B^{\rm reg}_{j,u_j}
          \le\frac{2C_{\rm ab}}{1-\varrho}\,\varrho^D.
\end{equation}
At $D=\lceil\vartheta_*m\rceil$, with fixed
$0<\vartheta_*<1/2$, the product of this complete positive charge
with $m^2$ converges to zero, uniformly in the original radius and
exact labels and in the prescribed levels.  Its pushforward by the
unchanged physical roof has at most the same total variation.
\end{theorem}
\begin{proof}
On the original endpoint cap, the identities in
Corollary~\ref{cor:bf-wronskian} give
\[
 f^\varepsilon
       =\frac{W^\varepsilon u_x}{4\pi c_*h_x},\qquad
 h_x=\frac{R^2\chi}{\det G}
       =R^2e^{\mathsf T}G^{-1}e,\qquad e=g/|g|.
\]
The last equality is the two-dimensional cofactor identity
$e^{\perp\mathsf T}Ge^\perp/\det G=e^{\mathsf T}G^{-1}e$.
Since $G\preceq\Lambda_\circ I$,
\begin{equation}\label{eq:ab-density-upper}
              f^\varepsilon
             \le\frac{\Lambda_\circ}{4\pi c_*R_-^2}\,u_x.
\end{equation}
The original bad-source weight satisfies $0\le W^\varepsilon\le1$;
it is not set to zero on a zero-clearance boundary.  There is exactly
one original conditioning factor in \eqref{eq:ab-density-upper}.

Represent the fixed-level two-sign angular source by the actual
one-sided interval partitions and their overlaps of
Theorem~\ref{thm:mb-two-sided-sum}.  On each overlap take the
supremum of the genuine root weight $u_x$, not a supremum on a
formal nonphysical extension.  Corollary~\ref{cor:ab-complete-piece-turning}
bounds the integral of $|\partial_\vartheta\theta|$ on its punctured
cap by $2\pi+V_A$.  Therefore
\[
 \mathfrak B^{\rm reg}_{j,u_0}
 \le\frac{\Lambda_\circ(2\pi+V_A)}{4\pi c_*R_-^2}
          \sum_{\text{actual overlaps}}\sup u_x.
\]
Insert the complete interval-overlay sum
\eqref{eq:mb-two-sided-sum}.  It gives exactly
\eqref{eq:ab-all-offset} and \eqref{eq:ab-final-constant}.
Countable partitions are handled first on finite compact regular
subfamilies, followed by nonnegative exhaustion.  Artificial seams
are assigned once on the angular source.  The overlap argument has
already paid its possible count; it is not paid again by a collision
word multiplier.  Original labels and valid one-sided Borel trace
selectors only restrict the positive measures at this stage.

This proof never divides the source into zero and nonzero formal
critical offsets.  In particular the term identified after
Theorem~\ref{thm:zf-flux} is retained through the exact turning
integral, even when it concentrates sharply.  The estimate bounds
that integral and does not assert a uniform pointwise normal speed.
Arbitrarily short genuine overlaps and components approaching a
critical record are included with the same constant.

Summing \eqref{eq:ab-all-offset} over $D\le j\le m-1-D$ gives
twice a subseries of $\sum_{j\ge D}\varrho^j$, proving
\eqref{eq:ab-full-depth}.  An empty index range contributes zero.
After the stated finite choice of $D$, multiplication by $m^2$
leaves a constant times $m^2\varrho^{\vartheta_*m}$, tending to zero.
All outer widths and the original central target remain fixed.
The supremum over radius, labels and levels is permitted by the
same displayed finite-count bound.  Pushforward cannot increase
the total mass of a positive measure, proving the roof-TV statement.
\end{proof}

\paragraph{The exact improvement and the remaining target.}
Theorem~\ref{thm:ab-all-offset} closes the unrestricted
\emph{regular marked-boundary} angular summation, including the
nonzero-offset term left in Section~\ref{sec:zf-zero-offset}.
The mechanism is the complete two-sided physical inverse-weight
sum and uniform total turning, rather than smallness of a local
offset.  It preserves every actual marked copy and every original
physical clock and record label.  No claim that a small total
variation is a small roof essential height is made.

Section~\ref{pc:section} additionally pays the leading incoming
zero-clearance marked critical fan when all visited incidences
remain positive, including collapsed slices and shallow marks.
The nonleading critical remainders, distinct visited-grazing and
section/pairing boundary measures, shallow low-load same-roof
shortage and whole physical comparison capacity retain their stated
obligations.  In particular \eqref{eq:ab-full-depth} does not imply
the original $m^2L^\infty$ local limit simply by multiplying a fixed
shallow-source mass bound by $m^2$.  The raw target
\eqref{eq:v76-target}, its exact arithmetic reference,
conditioning and original order of limits remain unchanged.
